% Part: normal-modal-logic
% Chapter: filtrations
% Section: finite

\documentclass[../../../include/open-logic-section]{subfiles}

\begin{document}

\olfileid[hi]{nml}{fil}{fin}

\olsection{निस्यंदन परिमित होते हैं}

हमने उप-!!{formula}ों के अधीन संवृत किसी भी समुच्चय~$\Gamma$ के लिए
निस्यंदन परिभाषित किए हैं। परिभाषा में स्वयं कुछ भी यह प्रत्याभूत नहीं
करता कि निस्यंदन परिमित हैं। वास्तव में जब $\Gamma$ अपरिमित हो (जैसे सभी
!!{formula}ों का समुच्चय), तो निस्यंदन भी अपरिमित हो सकता है। किंतु यदि
$\Gamma$ परिमित है (जैसे किसी दिए गए !!{formula}~$!A$ के उप-!!{formula}ों का
समुच्चय), तो $\Gamma$ के माध्यम से प्रत्येक निस्यंदन भी परिमित है।

\begin{prop}\ollabel{prop:filt-are-finite}
  यदि $\Gamma$ परिमित है, तो $\Gamma$ के माध्यम से किसी प्रतिरूप
  $\mModel{M}$ का प्रत्येक निस्यंदन $\mModel{M^*}$ भी परिमित है।
\end{prop}

\begin{proof}
  $W^*$ का आकार तुल्यता संबंध~$\equiv$ के अधीन भिन्न वर्गों~$[w]$ की
  संख्या है। ऐसे किसी वर्ग के कोई दो जगत $u$, $v$---अर्थात् ऐसे $u$ और
  $v$ कि $u \equiv v$---$\Gamma$ के सभी !!{formula}ों~$!A$ पर सहमत होते
  हैं: $!A \in \Gamma$ या तो $u$ और $v$, दोनों पर सत्य है, या किसी पर
  नहीं। अतः प्रत्येक वर्ग~$[w]$, $\Gamma$ के किसी उपसमुच्चय के अनुरूप है,
  अर्थात् उन सभी $!A \in \Gamma$ के समुच्चय के, जो $[w]$ के जगतों पर
  सत्य हैं। कोई दो भिन्न वर्ग $[u]$ और $[v]$, $\Gamma$ के एक ही उपसमुच्चय
  के अनुरूप नहीं हैं। क्योंकि यदि $u$ पर सत्य !!{formula}ों का समुच्चय और
  $v$ पर सत्य !!{formula}ों का समुच्चय समान हों, तो $u$ और $v$, $\Gamma$ के
  सभी सूत्रों पर सहमत हैं, अर्थात् $u \equiv v$। किंतु तब $[u] = [v]$।
  अतः $W^*$ से $\Pow{\Gamma}$ में कोई !!a{injective} फलन है, और इसलिए
  $\card{W^*} \le \card{\Pow{\Gamma}}$। अतः यदि $\Gamma$ में $n$ वाक्य
  हैं, तो $W^*$ की गणनीयता $2^n$ से अधिक नहीं है।
\end{proof}

\end{document}


